\documentclass{article}
\usepackage[T1]{fontenc}
\usepackage{lmodern}
\usepackage{graphicx}
\usepackage{amsmath,amssymb}
\usepackage[a4paper,margin=2cm]{geometry}
\usepackage[absolute,overlay]{textpos}
\setlength{\TPHorizModule}{1mm}
\setlength{\TPVertModule}{1mm}
\usepackage[indonesian]{babel}
\usepackage{hyperref}
\setlength{\emergencystretch}{2em}
% unit: Halaman 4

\begin{document}

% === Halaman 4 (llm) ===

\begin{itemize}
  \item Alternatif 2: menyembunyikannya di dalam tulisan lain
\end{itemize}

\textcolor{blue}{masihkah} \textcolor{blue}{ada} \textcolor{blue}{lara} \textcolor{blue}{apabila} \textcolor{blue}{memoriku} \textcolor{red}{ingat} \textcolor{blue}{nestapa} \textcolor{blue}{itu.} \textcolor{red}{kita} \textcolor{red}{ingin} \textcolor{red}{tetap}
\textcolor{blue}{abadikan} \textcolor{blue}{kisah} \textcolor{blue}{asmara.} \textcolor{red}{bersamamu} \textcolor{blue}{usiaku} \textcolor{red}{renta.}

\vspace{5mm}

Wendy tidak akan curiga!

\vspace{5mm}

Information hiding dengan steganografi!

\end{document}
